% Chapter 12 — SQL Patterns and Worked Examples

\chapter{SQL Patterns and Worked Examples}
\label{ch:sqlpatterns}

\section{Purpose}
These queries are templates for demos and vivas. Replace \texttt{\$1} with parameters in application code or Supabase filters.

\section{Roster listing}
Listing~\ref{lst:roster}: everyone enrolled in a course with readable status text.

\begin{lstlisting}[caption={Students enrolled in one course (with status label)},label={lst:roster}]
SELECT
  e.enrollment_id,
  u.user_id,
  u.full_name,
  u.email,
  es.status_name,
  e.enrolled_at
FROM enrollments e
JOIN users u
  ON u.user_id = e.user_id
JOIN enrollment_statuses es
  ON es.status_id = e.status_id
WHERE e.course_id = $1
ORDER BY e.enrolled_at DESC;
\end{lstlisting}

\section{Course-level summary (illustrative)}
Listing~\ref{lst:coursesummary}: joins across enrollments, progress paths, and quizzes. \textbf{Note:} real production queries sometimes split these metrics or use views to avoid double-counting when joins multiply rows---use this example to discuss join grain in class.

\begin{lstlisting}[caption={Illustrative per-course aggregates (discuss join grain carefully)},label={lst:coursesummary}]
SELECT
  c.course_id,
  c.title,
  COUNT(DISTINCT e.user_id) AS enrolled_students,
  COALESCE(ROUND(AVG(cp.progress_percent)::numeric, 2), 0) AS avg_progress_pct,
  COALESCE(ROUND(AVG(qa.score)::numeric, 2), 0) AS avg_quiz_score
FROM courses c
LEFT JOIN enrollments e
  ON e.course_id = c.course_id
LEFT JOIN content ct
  ON ct.course_id = c.course_id
LEFT JOIN content_progress cp
  ON cp.content_id = ct.content_id AND cp.user_id = e.user_id
LEFT JOIN quiz q
  ON q.course_id = c.course_id
LEFT JOIN quiz_attempts qa
  ON qa.quiz_id = q.quiz_id AND qa.user_id = e.user_id
WHERE c.deleted_at IS NULL
GROUP BY c.course_id, c.title;
\end{lstlisting}

\section{Fixing SERIAL sequences after imports}
If you bulk-insert explicit IDs, auto-increment sequences may fall behind the actual maximum ID. Listing~\ref{lst:setval} resets the sequence safely.

\begin{lstlisting}[caption={Align SERIAL sequence with current max ID},label={lst:setval}]
SELECT setval(
  pg_get_serial_sequence('public.courses', 'course_id'),
  COALESCE((SELECT MAX(course_id) FROM public.courses), 1),
  true
);
\end{lstlisting}

\section{Presentation tip}
Walk through one listing line by line: FROM/WHERE choose rows; JOIN connects tables; GROUP BY sets aggregation buckets.
